\documentclass[10pt,a4paper]{amsart}
\usepackage[margin=19mm]{geometry}
\usepackage{amsmath,amssymb,mathtools}
\usepackage[scr=rsfs,cal=euler]{mathalfa}
\usepackage{xfrac}
\usepackage[unicode,hidelinks,bookmarks=false]{hyperref}
\newcommand{\ie}{i.e.\ }
\newcommand{\cf}{cf.\ }
\newcommand{\resp}{respectively\ }
\newcommand{\Subsection}{\subsection}
\providecommand{\nobreakdash}{\nobreak\hbox{-}}
\setlength{\emergencystretch}{2em}
\multlinegap0pt
\usepackage{bm}
\usepackage{bbm}
\usepackage{dsfont}
\usepackage{xspace}
\usepackage{cancel}
\usepackage{mathtools}
\usepackage{amsthm}
\makeatletter
\@ifundefined{theorem}{\newtheorem{theorem}{Theorem}}{}
\@ifundefined{lem}{\newtheorem{lem}[theorem]{Lemma}}{}
\makeatother
\makeatletter
\expandafter\ifx\csname proof\endcsname\relax
\renewenvironment{proof}[1][\proofname]{\par
  \pushQED{\qed}%
  \normalfont \topsep6\p@\@plus6\p@\relax
  \trivlist
  \item[\hskip\labelsep\bfseries #1\@addpunct{.}]%
}{%
  \popQED\endtrivlist\@endpefalse
}
\fi
\makeatother
\title[Forward and backward problems for abstract time-fractional S]{Forward and backward problems for abstract time-fractional Schr\"odinger equations}
\author{S. E. Chorfi, F. Et-tahri, L. Maniar, M. Yamamoto}
\date{Source: arXiv:2510.03600v1 (CC BY 4.0)}
\begin{document}
\maketitle

\noindent\emph{Source and licence.} Forward and backward problems for abstract time-fractional Schr\"odinger equations. Version arXiv:2510.03600v1 (CC BY 4.0), distributed under the Creative Commons Attribution 4.0 International licence, as recorded in the arXiv licence metadata for this exact version and shown on the abstract page https://arxiv.org/abs/2510.03600v1. Full rights notice: licenses/arxiv-2510.03600-CC-BY-4.0.txt.\par\smallskip

\noindent\textit{Editorial scope.} Counted unit: the Theorem whose source label is \texttt{None} in the article (source0.tex lines 983--993, source label \texttt{None}), delivered together with the article's own complete original proof of it (source0.tex lines 994--1053, one \texttt{proof} environment of 60 lines). No mathematics is written, repaired or re-derived: statement and proof are the article's own text line for line. Required context retained uncounted: es1 (lines 193--198); es3 (lines 193--198); solution formula3 (lines 797--799); D2 (lines 802--806); fsweq (lines 906--911); asypsi (lines 917--922); asyE (lines 926--932). Every definition, display and label of those lines is retained unchanged. Omitted from this package and not redistributed: 1--192, 199--796, 800--801, 807--905, 912--916, 923--925, 933--982, 1054--1229. Nothing inside a retained excerpt is omitted or altered: every retained body is delivered exactly as the article's lines are written, so no omission receipt of any kind accompanies this package. Citations: the retained excerpts carry no citation command at all, so no bibliography entry of the article is transcribed and no reference is redistributed. Reference adaptation: none was needed and none was made; every reference inside the delivered unit resolves inside it, so no unresolved reference or citation is produced. Printed slips kept literal: no printed word, symbol, hypothesis or number of the counted statement or of its proof was altered. Layout and class: the article's own class is \texttt{amsart} with packages including amssymb, amsmath, amsthm, amscd, mathrsfs, paralist, color, cite, bigints, dsfont, empheq, cases, enumitem, caption; the wrapper is the lane's pinned \texttt{amsart} editorial wrapper at 10pt on A4, which restates the theorem-like environments the retained text needs and adds the article's own macro definitions that the retained text needs (none), with extra wrapper packages (bm, bbm, dsfont, xspace, cancel, mathtools). The delivered numbering is the unit's own. Assets: no retained span contains a figure environment, a graphics-inclusion command, a PDF-page inclusion, a table or a TikZ picture; no image file is copied into this package, the package declares no asset dependency and no image file is redistributed with it. Licence: version arXiv:2510.03600v1 of this article is distributed under the Creative Commons Attribution 4.0 International licence, as recorded in the arXiv licence metadata for this exact version and shown on the abstract page https://arxiv.org/abs/2510.03600v1. A full rights notice is delivered at licenses/arxiv-2510.03600-CC-BY-4.0.txt. Characters: every retained excerpt is pure ASCII, so the delivered engine, XeLaTeX with its default Latin Modern text font, reports no missing character.
\begin{align}
& \left|E_{\alpha, 1}\left(\mathrm{i}^{-\alpha} \lambda t^\alpha\right)\right| \leq C_1+\frac{C_2}{1+\lambda t^\alpha} \leq C_3; \label{es1}\\
& \left|E_{\alpha, 2}\left(\mathrm{i}^{-\alpha} \lambda t^\alpha\right)\right| \leq C_1\left(1+\lambda t^\alpha\right)^{\frac{-1}{\alpha}}+\frac{C_2}{1+\lambda t^\alpha} \leq \frac{C_3}{\left(1+\lambda t^\alpha\right)^{\min(1,\frac{1}{\alpha})}}; \label{es2}\\
& \left|E_{\alpha, \alpha}\left(\mathrm{i}^{-\alpha} \lambda t^\alpha\right)\right| \leq C_1\left(1+\lambda t^\alpha\right)^{\frac{1-\alpha}{\alpha}}+\frac{C_2}{1+\lambda t^\alpha} \leq C_3\left(1+\lambda t^\alpha\right)^{\frac{1-\alpha}{\alpha}}; \label{es3}\\
& \left|E_{\alpha, \alpha-1}\left(\mathrm{i}^{-\alpha} \lambda t^\alpha\right)\right| \leq C_1\left(1+\lambda t^\alpha\right)^{\frac{2-\alpha}{\alpha}}+\frac{C_2}{1+\lambda t^\alpha} \leq C_3\left(1+\lambda t^\alpha\right)^{\frac{2-\alpha}{\alpha}}. \label{es4}
\end{align}

	\begin{align}\label{solution formula3}
        u(t)=\sum_{n=1}^{\infty}\left\langle u_0, \varphi_n\right\rangle E_{\alpha, 1}\left(\mathrm{i}^{-\alpha}\lambda_n t^\alpha\right) \varphi_n + \sum_{n=1}^{\infty}\left\langle u_1, \varphi_n\right\rangle tE_{\alpha, 2}\left(\mathrm{i}^{-\alpha}\lambda_n t^\alpha\right) \varphi_n,
	\end{align}

\begin{align}
	\partial_t u(t)=&\sum_{n=1}^{\infty}\mathrm{i}^{-\alpha}\lambda_n t^{\alpha-1}\left\langle u_0, \varphi_n\right\rangle E_{\alpha, \alpha}\left(\mathrm{i}^{-\alpha}\lambda_n t^\alpha\right) \varphi_n+ \sum_{n=1}^{\infty}\left\langle u_1, \varphi_n\right\rangle E_{\alpha, 1}\left(\mathrm{i}^{-\alpha}\lambda_n t^\alpha\right) \varphi_n, \label{D2}\\
	\partial_t^2 u(t)=&\sum_{n=1}^{\infty}\mathrm{i}^{-\alpha}\lambda_n t^{\alpha-2}\left\langle u_0, \varphi_n\right\rangle E_{\alpha, \alpha-1}\left(\mathrm{i}^{-\alpha}\lambda_n t^\alpha\right) \varphi_n \nonumber\\
    &+\sum_{n=1}^{\infty}\mathrm{i}^{-\alpha}\lambda_n t^{\alpha-1}\left\langle u_1, \varphi_n\right\rangle E_{\alpha, \alpha}\left(\mathrm{i}^{-\alpha}\lambda_n t^\alpha\right) \varphi_n. \label{D3}
\end{align}

\begin{equation}\label{fsweq}
\begin{cases}
\mathrm{i}^{\alpha}\partial_t^\alpha u +Au=0, &\qquad t \in (0,T),\\
u(0)=u_0, \qquad \partial_t u(0)=u_1.
\end{cases}
\end{equation}

\begin{lem}
The function $\psi(t)$ satisfies the following asymptotic expansion
\begin{equation}\label{asypsi}
\psi(t)=\frac{\mathrm{i}^{\alpha-1}}{\alpha\Gamma(2-\alpha)} t^{\frac{1}{\alpha}-1}e^{-\mathrm{i}t^{\frac{1}{\alpha}}}-\frac{2 \mathrm{i}^\alpha}{\alpha\Gamma(1-\alpha)t}e^{-\mathrm{i}t^{\frac{1}{\alpha}}}+\mathcal{O}\left(\frac{1}{t^{1+\frac{1}{\alpha}}}\right), \; t\to \infty.
\end{equation}
\end{lem}

\begin{equation}\label{asyE}
\begin{aligned}
E_{\alpha, 1}(\mathrm{i}^{-\alpha}t)&=\frac{1}{\alpha} e^{-\mathrm{i}t^{\frac{1}{\alpha}}}-\frac{1}{\Gamma(1-\alpha) \mathrm{i}^{-\alpha}t}+\mathcal{O}\left(\frac{1}{t^2}\right), \\
E_{\alpha, 2}(\mathrm{i}^{-\alpha}t)&=\frac{1}{\alpha} \mathrm{i}t^{\frac{-1}{\alpha}} e^{-\mathrm{i}t^{\frac{1}{\alpha}}}-\frac{1}{\Gamma(2-\alpha) \mathrm{i}^{-\alpha}t}+\mathcal{O}\left(\frac{1}{t^2}\right), \\
tE_{\alpha, \alpha}(\mathrm{i}^{-\alpha}t)&=\frac{1}{\alpha} \mathrm{i}^{\alpha-1}t^{\frac{1}{\alpha}} e^{-\mathrm{i}t^{\frac{1}{\alpha}}}-\frac{1}{\Gamma(-\alpha) \mathrm{i}^{-2\alpha}t} +\mathcal{O}\left(\frac{1}{t^2}\right), \; t\to \infty.
\end{aligned}
\end{equation}

\begin{theorem}
Assume that $T\notin \Lambda$. For any $u_T, v_T \in H_1$, there exist $u_0, u_1 \in H$ such that the solution $u$ to \eqref{fsweq} satisfies
$$
u(T)=u_T, \qquad \partial_t u(T)=v_T.
$$
Moreover, there exists a constant $C>0$ such that
\begin{equation}
\|u(0)\|+\|\partial_t u(0)\| \leq C\left(\left\|u(T)\right\|_{1}+\left\|\partial_t u(T)\right\|_{1}\right)
\end{equation}
for all $u_T, v_T \in H_1$.
\end{theorem}

\begin{proof}
We seek initial data $u_0, u_1 \in H$ such that the corresponding solution $u$ to \eqref{fsweq} satisfies $u(T)=u_T$ and $\partial_t u(T)=v_T.$ To do this, it is enough to find the coordinates of $u_0$ and $u_1$:
$$
u_{0n}=\langle u_0, \varphi_{n}\rangle, \quad u_{1n}=\langle u_1, \varphi_{n}\rangle.
$$
According to formulas \eqref{solution formula3} and \eqref{D2}, we have:
\begin{equation}\label{eq0}
\begin{aligned}
\|u(T)\|_{1}^2+\left\|\partial_t u(T)\right\|_{1}^2
%&= \sum_{n=1}^{\infty}  \lambda_n^2\left(\left|u_{0n} E_{\alpha, 1}\left(\mathrm{i}^{-\alpha}\lambda_n T^\alpha\right)+u_{1n} T E_{\alpha, 2}\left(\mathrm{i}^{-\alpha}\lambda_n T^\alpha\right)\right|^2\right. \\
%& \left.+\left|\mathrm{i}^{-\alpha}\lambda_n T^{\alpha-1} u_{0n} E_{\alpha, \alpha}\left(\mathrm{i}^{-\alpha}\lambda_n T^\alpha\right)+u_{1n} E_{\alpha, 1}\left(\mathrm{i}^{-\alpha}\lambda_n T^\alpha\right)\right|^2\right) \\
= \sum_{n=1}^{\infty}  \lambda_n^2\left(\left|a_n\right|^2+\left|b_n\right|^2\right),
\end{aligned}
\end{equation}
where 
\begin{equation}\label{sys}
\left\{\begin{array}{l}
a_{n}:=u_{0n} E_{\alpha, 1}\left(\mathrm{i}^{-\alpha}\lambda_n T^\alpha\right)+u_{1n} T E_{\alpha, 2}\left(\mathrm{i}^{-\alpha}\lambda_n T^\alpha\right), \\
b_{n}:=\mathrm{i}^{-\alpha}\lambda_n T^{\alpha-1} u_{0n} E_{\alpha, \alpha}\left(\mathrm{i}^{-\alpha}\lambda_n T^\alpha\right)+u_{1n} E_{\alpha, 1}\left(\mathrm{i}^{-\alpha}\lambda_n T^\alpha\right).
\end{array}\right.
\end{equation}
Since $\psi\left(\lambda_n T^\alpha\right) \neq 0$ for all $n \in \mathbb{N}$ ($T\notin \Lambda$), we can solve \eqref{sys} with respect to $u_{0n}$ and $u_{1n}$ (determinant is nonzero):
$$
\left\{\begin{array}{l}
u_{0n}=\frac{1}{\psi\left(\lambda_n T^\alpha\right)}\left(a_n E_{\alpha, 1}\left(\mathrm{i}^{-\alpha}\lambda_n T^\alpha\right)-b_n T E_{\alpha, 2}\left(\mathrm{i}^{-\alpha}\lambda_n T^\alpha\right)\right) \\
u_{1n}=\frac{1}{\psi\left(\lambda_n T^\alpha\right)}\left(a_n \mathrm{i}^{-\alpha}\lambda_n T^{\alpha-1} E_{\alpha, \alpha}\left(\mathrm{i}^{-\alpha}\lambda_n T^\alpha\right)+b_n E_{\alpha, 1}\left(\mathrm{i}^{-\alpha}\lambda_n T^\alpha\right)\right).
\end{array}\right.
$$
By \eqref{asyE} and \eqref{asypsi}, we can choose large constants $M>0$ such that
$$
\begin{gathered}
\left|E_{\alpha, 1}(\mathrm{i}^{-\alpha}t)\right| \leq \frac{1}{\alpha}+\frac{M_1}{t}, \quad\left|E_{\alpha, 2}(\mathrm{i}^{-\alpha}t)\right| \leq \frac{1}{\alpha} t^{-\frac{1}{\alpha}} + \frac{M_2}{t}, \\
\left|E_{\alpha, \alpha}(\mathrm{i}^{-\alpha}t)\right| \leq \frac{1}{\alpha} t^{\frac{1}{\alpha}-1} + \frac{M_3}{t^2}, \quad|\psi(t)| \geq \frac{1}{2\alpha \Gamma(2-\alpha)} t^{\frac{1}{\alpha}-1}, \quad t \geq M.
\end{gathered}
$$
Note that $\Gamma(2-\alpha)>0$. By fixing $N_0 \in \mathbb{N}$ and using \eqref{es1}-\eqref{es3},
\begin{align}
&\left|\psi\left(\lambda_n T^\alpha\right)\right| \geq \frac{1}{2 T^{\alpha-1} \Gamma(2-\alpha)} \frac{1}{\lambda_n^{1-\frac{1}{\alpha}}}, \quad n \geq N_0, \quad \left|E_{\alpha, 1}\left(\mathrm{i}^{-\alpha}\lambda_n T^\alpha\right)\right| \le C,\nonumber\\%\ge\frac{C}{\lambda_n^{2}}, \\
&\left|E_{\alpha, 2}\left(\mathrm{i}^{-\alpha}\lambda_n T^\alpha\right)\right| \le \frac{C}{\lambda_n^{\frac{1}{\alpha}}}, \quad\left|\lambda_n E_{\alpha, \alpha}\left(\mathrm{i}^{-\alpha}\lambda_n T^\alpha\right)\right| \leq C \lambda_n^{\frac{1}{\alpha}}, \quad n\in \mathbb{N}. \label{estEa}
\end{align}
Therefore,
\begin{equation}\label{esu0u1}
\left|u_{0n}\right| \leq C \lambda_n^{1-\frac{1}{\alpha}}\left( \left|a_n\right|+\left|b_n\right|\right), \quad\left|u_{1n}\right| \leq C \lambda_n\left(\left|a_n\right|+\left|b_n\right|\right), \quad n \geq N_0.
\end{equation}
Since $\psi\left(\lambda_n T^\alpha\right) \neq 0$ for each $n \in \mathbb{N}$, we have
$$
\begin{aligned}
& \left|u_{0n}\right| \leq C \left(\left|a_n\right|+\left|b_n\right|\right)\max_{1 \leq n \leq N_0-1}\left|\frac{1}{\psi\left(\lambda_n T^\alpha\right)}\right|, \\
& \left|u_{1n}\right| \leq C \left(\lambda_n^{\frac{1}{\alpha}}\left|a_n\right|+\left|b_n\right|\right)\max_{1 \leq n \leq N_0-1}\left|\frac{1}{\psi\left(\lambda_n T^\alpha\right)}\right|, \qquad 1 \leq n \leq N_0-1,
\end{aligned}
$$
Then \eqref{esu0u1} is satisfied for all $n \in \mathbb{N}$. Hence
$$
\sum_{n=1}^{\infty} \left(\left|u_{0n}\right|^2+\left|u_{1n}\right|^2\right) \leq C \sum_{n=1}^{\infty}  \lambda_n^2\left(\left|a_n\right|^2+\left|b_n\right|^2\right).
$$
Using \eqref{eq0}, we obtain
$$
\|u_0\|^2+\|u_1\|^2 \leq C \left(\|u(T)\|_{1}^2+\left\|\partial_t u(T)\right\|_{1}^2\right) .
$$
\end{proof}

\end{document}
